\section*{Заключение}
\addcontentsline{toc}{section}{Заключение}

В ходе выполнения курсовой работы была разработана экспертная система для выбора параметров шрифтовой гарнитуры при подготовке бейджей к офсетной печати. Была описана предметная область, выделены основные критерии выбора, построено бинарное дерево решений и выполнена реализация модели в среде CLIPS. Дерево содержит 15 правил, 16 фактов и 6 ярусов. Работа системы продемонстрирована на нескольких сценариях, в которых программа задаёт пользователю вопросы и выдаёт рекомендации.

\subsection*{Достоинства реализации}

\begin{itemize}
    \item База знаний опирается на рекомендации ISO~24509:2019 по читаемости текста и охватывает наиболее частые сценарии подготовки бейджей.
    \item Сложность алгоритма составляет $O(\log n)$, где $n$~--- число конечных решений.
\end{itemize}

\subsection*{Недостатки реализации}

\begin{itemize}
    \item При добавлении новых правил может потребоваться перестроение поддерева, корнем которого является изменяемый узел-правило.
    \item Система не допускает неопределённых ответов пользователя (напимер, <<не знаю>>), поэтому для получения результата необходимо дать однозначный ответ на каждый задаваемый вопрос.
\end{itemize}

\subsection*{Масштабируемость}

\begin{itemize}
    \item Добавление графического интерфейса.
    \item Упрощение процесса добавление новых знаний, например, через графический интерфейс.
\end{itemize}

Таким образом, разработанная экспертная система выполняет поставленную задачу и может быть расширена для более детальной типографической консультации.
